\usemodule[memos][fontstyle=ss]

\starttext

\chapter{startDualText 合并组件测试}

% ============================================================
\section{基础功能测试}

\subsection{完整双栏（默认）}

\startDualText[number={No.001},title={这是标题}]
这是一个完整的标题条+正文示例。左侧编号，右侧标题，下方正文。
\stopDualText
% 【预期】顶部标题条：左编号底色块 + 右标题；下方正文框完整；接缝单线贴合。

\subsection{仅显示编号}

\startDualText[number={A-123},showtitle=0]
只显示编号，不显示标题。
\stopDualText
% 【预期】标题条只有左侧编号块，右侧无标题框；正文正常。

\subsection{仅显示标题}

\startDualText[title={仅标题示例},shownumber=0]
只显示标题，不显示编号。
\stopDualText
% 【预期】标题条只有标题块（占满或按 titlepos）；正文正常。

\subsection{无编号和标题}

\startDualText[shownumber=0,showtitle=0]
无标题条，纯正文框。
\stopDualText
% 【预期】无标题条，不占标题高度；纯正文框。

% ============================================================
\section{标题位置 titlepos 测试}

\subsection{titlepos=1 比例占满（默认）}

\startDualText[number={No.1},title={比例占满标题},titlepos=1]
titlepos=1 时编号和标题按 numberratio:titleratio 占满全宽。
\stopDualText
% 【预期】编号块与标题块各占约一半宽度，无空隙。

\subsection{titlepos=0 按文字宽 fit}

\startDualText[number={No.2},title={按文字宽度},titlepos=0]
titlepos=0 时按文字实际宽度排布。
\stopDualText
% 【预期】编号块和标题块宽度贴合文字，两侧留白。

\subsection{titlepos=0 + 左对齐}

\startDualText[number={No.3},title={左对齐},titlepos=0,titlealign=1]
标题左对齐。
\stopDualText
% 【预期】fit 盒子靠左。

\subsection{titlepos=0 + 右对齐}

\startDualText[number={No.4},title={右对齐},titlepos=0,titlealign=3]
标题右对齐。
\stopDualText
% 【预期】fit 盒子靠右。

\subsection{titlepos=1 + 编号左/标题右}

\startDualText[number=01,title=标题右对齐,titlepos=1,numberalign=1,titlealign=3,
  numberframe=0,titleframe=0,leftframe=0,rightframe=0,bottomframe=0,topframe=1]
编号左对齐，标题右对齐，仅上边框。
\stopDualText
% 【预期】仅一条上边框线；编号文字靠左、标题文字靠右。

\subsection{titlepos=1 + 仅编号右对齐}

\startDualText[number=04,showtitle=0,titlepos=1,numberalign=3,
  numberframe=0,titleframe=0,leftframe=0,rightframe=0,bottomframe=0,topframe=1]
仅编号时右对齐。
\stopDualText
% 【预期】编号文字靠右，仅上边框。

% ============================================================
\section{标题栏边框控制}

\subsection{隐藏 number 边框}

\startDualText[number={No.001},title={隐藏 number 边框},numberframe=0]
只隐藏编号块的边框。
\stopDualText
% 【预期】编号块有底色无边框；标题块有边框。

\subsection{隐藏 title 边框}

\startDualText[number={No.002},title={隐藏 title 边框},titleframe=0]
只隐藏标题块的边框。
\stopDualText
% 【预期】标题块有底色无边框；编号块有边框。

\subsection{隐藏所有标题栏边框}

\startDualText[number={No.003},title={隐藏所有边框},numberframe=0,titleframe=0]
编号和标题块都无边框。
\stopDualText
% 【预期】两块都只有底色，无边框线。

\subsection{自定义边框颜色}

\startDualText[number={No.004},title={自定义边框颜色},
  numberframecolor=red,titleframecolor=blue]
编号边框红色，标题边框蓝色。
\stopDualText
% 【预期】编号块红框、标题块蓝框。

% ============================================================
\section{宽度比例测试}

\subsection{1:2 比例}

\startDualText[number={No.},title={标题宽度是编号的两倍},
  numberratio=1,titleratio=2,titlepos=1]
标题宽度是编号的两倍。
\stopDualText
% 【预期】标题块宽度约为编号块两倍。

\subsection{1:3 比例}

\startDualText[number={A},title={标题宽度是编号的三倍},
  numberratio=1,titleratio=3,titlepos=1]
标题宽度是编号的三倍。
\stopDualText
% 【预期】标题块宽度约为编号块三倍。

% ============================================================
\section{圆角测试}

\subsection{标题条圆角 hdrradius}

\startDualText[number={圆角},title={标题条圆角},hdrradius=6pt]
hdrradius 设置标题块圆角。
\stopDualText
% 【预期】标题块外侧圆角；接缝处保持直角。

\subsection{正文圆角 radius}

\startDualText[number={R},title={正文圆角},radius=6pt]
radius 设置正文框圆角。
\stopDualText
% 【预期】正文框四角圆角；标题块直角。

\subsection{正文圆角 + bodycorner 选角}

\startDualText[number={C},title={仅下两角圆},radius=6pt,bodycorner=2]
bodycorner=2 只圆下两角。
\stopDualText
% 【预期】正文框仅下方两角圆角，上方两角直角。

\subsection{标题圆角 + 正文圆角组合}

\startDualText[number={组合},title={标题+正文圆角},
  hdrradius=5pt,radius=8pt]
标题块与正文框各自圆角。
\stopDualText
% 【预期】标题块圆角、正文框圆角；接缝处各自独立。

% ============================================================
\section{背景样式测试}

\subsection{simple（默认）}

\startDualText[number={S},title={simple 样式},]
默认简单样式，纯色背景+边框。
\stopDualText
% 【预期】纯色背景，完整边框。

\subsection{shadowbox}

\startDualText[number={Sh},title={阴影框},
  shadow=yes,
  shadowcolor=lightthemecolor]
阴影框样式。
\stopDualText
% 【预期】右下角阴影效果。

\subsection{charbox}

\startDualText[number={Ch},title={字符图案},character={$\star$}]
字符水印样式。
\stopDualText
% 【预期】背景有旋转 45° 字符水印。

\subsection{framebox}

\startDualText[number={F},title={棋盘格},framebox=yes,
  framecolor=themecolor]
棋盘格图案样式。
\stopDualText
% 【预期】背景有棋盘格点阵图案。

\subsection{barbox}

\startDualText[number={B},title={左侧色条},bar=yes,
  framecolor=themecolor]
左侧色条样式。
\stopDualText
% 【预期】正文左侧有色条 + 三边框。

\subsection{symbolbox}

\startDualText[number={Y},title={角标符号},symbol=※]
右下角角标样式。
\stopDualText
% 【预期】正文右下角有角标符号。

\subsection{自定义背景色}

\startDualText[number={Bg},title={自定义底色},
  backgroundcolor=softthemecolor,
  numberbgcolor=orange,numbercolor=white,
  titlecolor=red,titlebgcolor=white,
  ]
自定义编号/标题/正文底色。
\stopDualText
% 【预期】编号块橙底白字；标题块白底红字；正文浅色底。

% ============================================================
\section{边框控制测试}

\subsection{仅左边框}

\startDualText[shownumber=0,showtitle=0,
  leftframe=1,rightframe=0,topframe=0,bottomframe=0]
只显示左边框。
\stopDualText
% 【预期】仅左侧一条竖线。

\subsection{仅左右边框}

\startDualText[shownumber=0,showtitle=0,
  leftframe=1,rightframe=1,topframe=0,bottomframe=0]
只显示左右边框。
\stopDualText
% 【预期】左右两条竖线，上下无线。

\subsection{仅上边框}

\startDualText[shownumber=0,showtitle=0,
  leftframe=0,rightframe=0,topframe=1,bottomframe=0]
只显示上边框。
\stopDualText
% 【预期】仅顶部一条横线。

\subsection{无边框}

\startDualText[shownumber=0,showtitle=0,
  leftframe=0,rightframe=0,topframe=0,bottomframe=0,
  backgroundcolor=softthemecolor]
不显示任何边框，仅底色。
\stopDualText
% 【预期】无边框线，仅可见底色区域。

\subsection{粗线 rulethickness}

\startDualText[number={粗},title={6pt 粗线},rulethickness=6pt]
粗线效果。
\stopDualText
% 【预期】边框线 6pt 粗；线外缘不跃出区域边界。

% ============================================================
\section{框线偏移测试}

\subsection{frameoffset=0（默认，线贴边界）}

\startDualText[number={F0},title={frameoffset 0},frameoffset=0pt,
  backgroundcolor=softthemecolor]
线外缘与背景边界齐平。
\stopDualText
% 【预期】边框线外缘与底色外缘齐平，不跃出。

\subsection{frameoffset=-2pt（官方居中跃出）}

\startDualText[number={F-},title={frameoffset -2pt},frameoffset=-2pt,
  backgroundcolor=softthemecolor]
线中心落在边界，左右各跃出 th/2。
\stopDualText
% 【预期】边框线向外跃出约 1pt（th/2）。

\subsection{frameoffset=+2pt（线更靠内）}

\startDualText[number={F+},title={frameoffset +2pt},frameoffset=2pt,
  backgroundcolor=softthemecolor]
线更靠内。
\stopDualText
% 【预期】边框线向内缩进，底色外缘露出一圈。

\subsection{backgroundoffset=0 vs 4pt}

\startDualText[number={B0},title={bgoff 0},backgroundoffset=0pt,
  backgroundcolor=softthemecolor]
背景不外扩。
\stopDualText

\startDualText[number={B4},title={bgoff 4pt},backgroundoffset=4pt,
  backgroundcolor=softthemecolor]
背景向外扩 4pt。
\stopDualText
% 【预期】第二例底色比第一例向外扩约 4pt，且不被裁切。

% ============================================================
\section{跨页测试 hdrrepeat}

\subsection{hdrrepeat=1 每页重复（默认）}

\startDualText[number={P1},title={每页标题},
  backgroundcolor=lightgray]
  \par 第一段。第二段。第三段。第四段。第五段。
  \par 第六段。第七段。第八段。第九段。第十段。
  \par 第十一段。第十二段。第十三段。第十四段。第十五段。
  \par 第十六段。第十七段。第十八段。第十九段。第二十段。
  \par 第二十一段。第二十二段。第二十三段。第二十四段。第二十五段。
  \par 第二十六段。第二十七段。第二十八段。第二十九段。第三十段。最后一段。
\stopDualText
% 【预期】每一页顶部都有标题条；各页正文框背景与框线完整。

\subsection{hdrrepeat=0 仅首页}

\startDualText[number={P2},title={仅首页标题},
  backgroundcolor=lightgray,hdrrepeat=0]
  \par 第一段。第二段。第三段。第四段。第五段。
  \par 第六段。第七段。第八段。第九段。第十段。
  \par 第十一段。第十二段。第十三段。第十四段。第十五段。
  \par 第十六段。第十七段。第十八段。第十九段。第二十段。
  \par 第二十一段。第二十二段。第二十三段。第二十四段。第二十五段。
  \par 第二十六段。第二十七段。第二十八段。第二十九段。第三十段。最后一段。
\stopDualText
% 【预期】仅首页顶部有标题条；后续页正文框完整（可能因 topoffset 偏高）。

% ============================================================
\section{组合效果测试}

\subsection{阴影 + 圆角}

\startDualText[number={组合},title={阴影+圆角},
  shadow=yes,
  bodycorner=2,
  hdrcorner=1,
  hdrradius=8pt,
  radius=8pt,shadowcolor=lightthemecolor]
阴影和圆角的组合。
\stopDualText
% 【预期】正文框圆角+阴影；标题条正常。

\subsection{左侧条 + 粗线}

\startDualText[number={组合},title={色条+粗线},
  bar=yes,rulethickness=2pt,leftframe=1]
左侧色条与粗边框组合。
\stopDualText
% 【预期】左侧色条 + 2pt 边框线。

\subsection{字符图案 + 圆角}

\startDualText[number={组合},title={字符+圆角},
  character={$\star$},
  character={$\heartsuit$},
  bodycorner=2,
  hdrcorner=1,
  hdrradius=10pt,
  radius=10pt]
字符水印与圆角组合。
\stopDualText
% 【预期】圆角正文框 + 爱心水印。

\subsection{仅左边框模拟 blockbox}

\startDualText[number={模拟},title={模拟 blockbox},
  leftframe=1,rightframe=0,topframe=0,bottomframe=0,rulethickness=2pt]
仅左边框模拟 blockquote 效果。
\stopDualText
% 【预期】仅左侧 2pt 竖线，模拟引用块。

\subsection{自定义样式}

\startDualText[number={? 问题探究},title={自定义样式},
  numberbgcolor=orange,numbercolor=white,
  titlecolor=red,titlebgcolor=white,
  backgroundcolor=lightorange,
  numberframe=0,titleframe=0,
  leftframe=0,rightframe=0,topframe=1,bottomframe=0,
  numberratio=1,titleratio=3]
自定义编号/标题/正文配色与边框。
\stopDualText
% 【预期】橙底白字编号 + 红字白底标题 + 浅橙正文 + 仅上边框；标题占 3/4 宽。

% ============================================================
\section{混合样式测试}

\subsection{shadow + bar}

\startDualText[number={X1},title={阴影+色条},
  shadow=yes,bar=yes,
  framecolor=themecolor,
  shadowcolor=lightthemecolor]
同时开启阴影与左色条。
\stopDualText
% 【预期】右下阴影 + 左色条 + 边框，叠加正常。

\subsection{shadow + character}

\startDualText[number={X2},title={阴影+字符},
  shadow=yes,character={$\heartsuit$},
  framecolor=red,
  shadowcolor=lightthemecolor]
同时开启阴影与字符水印。
\stopDualText
% 【预期】阴影 + 爱心水印，水印不被阴影遮盖。

\subsection{bar + character}

\startDualText[number={X3},title={色条+字符},
  bar=yes,character={$\star$},
  framecolor=themecolor]
同时开启左色条与字符水印。
\stopDualText
% 【预期】左色条 + 星形水印，两者互不干扰。

\subsection{bar + framebox}

\startDualText[number={X4},title={色条+棋盘},
  bar=yes,framebox=yes,
  framecolor=themecolor]
同时开启左色条与棋盘格点阵。
\stopDualText
% 【预期】左色条覆盖在棋盘格之上，底色仍为白。

\subsection{shadow + bar + character}

\startDualText[number={X5},title={阴影+色条+字符},
  shadow=yes,bar=yes,character={$\heartsuit$},
  framecolor=themecolor,rulethickness=1.5pt,
  shadowcolor=lightthemecolor]
三者同时开启。
\stopDualText
% 【预期】阴影 + 左色条 + 爱心水印共存。

\subsection{shadow + bar + framebox}

\startDualText[number={X6},title={阴影+色条+棋盘},
  shadow=yes,bar=yes,framebox=yes,
  framecolor=themecolor,
  shadowcolor=lightthemecolor]
阴影 + 色条 + 棋盘格。
\stopDualText
% 【预期】阴影 + 左色条 + 棋盘格，色条不被棋盘遮住。

\subsection{shadow + character + symbol}

\startDualText[number={X7},title={阴影+字符+角标},
  shadow=yes,character={$\star$},symbol=※,
  framecolor=themecolor,
  shadowcolor=lightthemecolor]
阴影 + 水印 + 右下角标。
\stopDualText
% 【预期】阴影 + 星形水印 + ※ 角标同时可见。

\subsection{全套组合}

\startDualText[number={X8},title={全套},
  shadow=yes,bar=yes,framebox=yes,
  character={$\heartsuit$},symbol=※,
  framecolor=themecolor,rulethickness=1.5pt,radius=6pt,
  bodycorner=2,
  shadowcolor=lightthemecolor]
全部装饰同时开启。
\stopDualText
% 【预期】阴影 + 色条 + 棋盘 + 水印 + 角标 + 圆角共存，层次不乱。

\stoptext
